On six lidar dune fields, a 30\,m elevation model hides most slip faces, and neither a physically motivated detector nor models trained on lidar recover enough of them to make planned routes safer; on the real Copernicus product the loss is larger still. At 10\,m every method sees most faces and the best hazard layer more than doubles mission success over the slip model alone, but only a face map at about 5\,m makes routes safe. Map resolution, not the hazard model, is the binding constraint for small rovers on dunes. The next step is to test the same benchmark on Arabian dune fields with a drone survey and on a real rover whose slip has been measured on the same sand.
